\documentclass[12pt,a4paper]{article}
\usepackage[utf8]{inputenc}
\usepackage[italian]{babel}
\usepackage{amsmath, amssymb, amsfonts}
\usepackage{graphicx}
\usepackage{pgfplots}
\pgfplotsset{compat=1.18}
\usepackage{geometry}
\geometry{margin=2.5cm}
\usepackage{caption}
\usepackage{hyperref}
\usepackage{titlesec}

\title{\centering Studio matematico sulla coscienza umana}
\author{Autore: Giuseppe Carlino}
\date{\today}

% Riduzione spaziatura sezioni
\titlespacing*{\section}{0pt}{1.2\baselineskip}{0.5\baselineskip}
\titlespacing*{\subsection}{0pt}{1\baselineskip}{0.3\baselineskip}

\captionsetup{justification=centering}

\begin{document}

\maketitle
\tableofcontents
\newpage

%%%%%%%%%%%%%%%%%%%%
\section{Legge 0 – Conoscenza}

\subsection{Formalizzazione concettuale}

\[
t_c \quad \text{tempo percepito}, \quad k \quad \text{conoscenza accumulata}
\]

\[
k = f_0(t_c), \quad f_0'>0
\]

Perturbazioni:

\[
\frac{dk}{dt_c} = f_0'(t_c) + \eta_k(t_c)
\]

\subsection{Interpretazione}

\begin{itemize}
    \item La conoscenza cresce con l’esperienza temporale.
    \item Eventi esterni o interni introducono fluttuazioni concettuali.
\end{itemize}

\subsection{Diagramma concettuale}
\begin{center}
\begin{tikzpicture}
\begin{axis}[
    width=14cm, height=8cm,
    axis lines=middle,
    xlabel={$t_c$},
    ylabel={$k$},
    grid=major,
]
\addplot[smooth, thick, blue] {0.5*x + 0.1*sin(deg(x))};
\node at (axis cs:4,2.5) {\scriptsize Conoscenza crescente con il tempo};
\end{axis}
\end{tikzpicture}
\end{center}

\newpage
%%%%%%%%%%%%%%%%%%%%
\section{Legge 1 – Tempo}

\subsection{Formalizzazione}

\[
\frac{dt_c}{d\tau} = 1 + \eta_t(\tau,F)
\]

\subsection{Interpretazione}

\begin{itemize}
    \item Il tempo percepito non coincide necessariamente con il tempo lineare.
    \item Influenza la conoscenza e quindi le scelte.
\end{itemize}

\subsection{Diagramma concettuale}
\begin{center}
\begin{tikzpicture}
\begin{axis}[
    width=14cm, height=8cm,
    axis lines=middle,
    xlabel={Tempo reale $\tau$},
    ylabel={Tempo percepito $t_c$},
    grid=major,
]
\addplot[smooth, thick, red] {x + 0.3*sin(deg(0.5*x))};
\node at (axis cs:4,4.5) {\scriptsize Fluttuazioni temporali soggettive};
\end{axis}
\end{tikzpicture}
\end{center}

\newpage
%%%%%%%%%%%%%%%%%%%%
\section{Legge 2 – Scelta}

\subsection{Formalizzazione}

\[
\frac{ds}{d\tau} = -\nu(s - s_e(k,t_c)) + \eta_s(\tau)
\]

\subsection{Interpretazione}

\begin{itemize}
    \item Le scelte dipendono dalla conoscenza accumulata e dal tempo percepito.
    \item La retroazione dalla volontà sarà introdotta nella Legge 3.
\end{itemize}

\subsection{Diagramma concettuale}
\begin{center}
\begin{tikzpicture}
\begin{axis}[
    width=14cm, height=8cm,
    axis lines=middle,
    xlabel={$k$},
    ylabel={$s$},
    grid=major,
]
\addplot[smooth, thick, green] coordinates {(0,0.1) (1,0.5) (2,1.0) (3,1.3) (4,1.8)};
\node at (axis cs:3,1.5) {\scriptsize Scelta in funzione di conoscenza e tempo};
\end{axis}
\end{tikzpicture}
\end{center}

\newpage
%%%%%%%%%%%%%%%%%%%%
\section{Legge 3 – Volontà}

\subsection{Formalizzazione}

\[
\frac{dw}{d\tau} = -\xi(w - w_e(t_c,k,s,F)) + \eta_w(\tau)
\]

\subsection{Interpretazione}

\begin{itemize}
    \item La volontà integra informazioni da tempo, conoscenza e scelta.
    \item Introduce retroazioni sulle scelte future.
\end{itemize}

\subsection{Diagramma concettuale}
\begin{center}
\begin{tikzpicture}
\begin{axis}[
    width=14cm, height=8cm,
    axis lines=middle,
    xlabel={$s$},
    ylabel={$w$},
    grid=major,
]
\addplot[smooth, thick, purple] coordinates {(0,0.1) (1,0.4) (2,0.9) (3,1.5) (4,2.0)};
\node at (axis cs:3,1.7) {\scriptsize Volontà in funzione di scelte, tempo e conoscenza};
\end{axis}
\end{tikzpicture}
\end{center}

\newpage
%%%%%%%%%%%%%%%%%%%%
\section{Legge 4 – Legge Generale della Coscienza Quadristato}

\subsection{Formalizzazione matematica}

\[
\mathbf{X} = 
\begin{pmatrix}
t_c \\ k \\ s \\ w
\end{pmatrix}, \quad
O = 
\begin{pmatrix}
t_e \\ k_e \\ s_e \\ w_e(F)
\end{pmatrix}, \quad
\kappa = \|\mathbf{X}-O\|
\]

\[
\frac{d\mathbf{X}}{d\tau} = -\Lambda(\mathbf{X}-O) + \mathbf{\eta}(\tau), \quad
\Lambda = \mathrm{diag}(\lambda, \mu, \nu, \xi)
\]

\subsection{Diagramma concettuale quadridimensionale}

\begin{center}
\begin{tikzpicture}
\begin{axis}[
    view={70}{25},
    width=14cm, height=10cm,
    axis lines=center,
    grid=major,
    xmin=0, xmax=6, ymin=0, ymax=4, zmin=0, zmax=5,
    axis line style={->}, ticks=none,
]
% Traiettoria concettuale
\addplot3[smooth, thick, blue] coordinates {
(0,0.5,0.2) (1,1,1) (2,1.5,2) (3,2,2.5) (4,2.5,3.2) (5,3,4)
};
% Punto equilibrio assoluto O
\addplot3[only marks, red, mark=*] coordinates {(3,2,2.5)};
\node at (axis cs:2.2,2.5,4.3) {\scriptsize $O=(t_e,k_e,s_e,w_e)$};

% Nuvola concettuale attorno a O
\addplot3[only marks, mark=*, mark size=0.7pt, gray, opacity=0.3] coordinates {
(3.05,2.05,2.55) (2.95,1.95,2.45) (3.1,2.0,2.6) (2.9,2.0,2.4)
};

\end{axis}
\end{tikzpicture}
\end{center}

\subsection{Interpretazione e implicazioni}

\begin{itemize}
    \item La coscienza umana è quadridimensionale: tempo, conoscenza, scelta e volontà interagiscono costantemente.
    \item $\kappa$ misura la distanza dall’equilibrio $O$.
    \item Le perturbazioni rendono la coscienza intrinsecamente imprevedibile.
    \item Sistemi artificiali possono simulare parte dei processi, ma non replicano la complessità integrata dei quattro stati.
    \item La \textbf{Teoria della Coscienza Quadristato} deriva direttamente dalla Legge Generale.
\end{itemize}

\end{document}